\section{1901 | As primeiras matrículas portuguesas}
\subsection*{Licenças de circulação de automóveis}

As primeiras matrículas surgiram a 3 de outubro de 1901, então denominadas
“licenças de circulação de automóveis”. Nesta altura eram emitidas pelos
Governos Civis dos Distritos e apenas era obrigatória uma chapa de matrícula na
parte de trás do veículo, embora as fotografias da época mostrem que por vezes
também era colocada à frente.

As combinações eram compostas pelo nome ou pela abreviatura da sede de distrito
que atribuiu a licença e pelo número, podendo assim ver-se matrículas como
\plate{LISBOA 130} ou \plate{LXA 256}.

O decreto com essa data estabelecia a obrigatoriedade de todos os automóveis
trazerem na parte de trás uma chapa metálica com os caracteres do número da
licença e da sede de distrito onde foi concedida.

O primeiro registo automóvel em Portugal foi feito em nome de Augusto Sabido
Teixeira de Aragão, de Santiago do Cacém.

\section{1911 | Série X-000}
\subsection*{1.ª padronização do sistema de matrículas português}

Com o aumento do número de carros em Portugal, o sistema foi padronizado a
nível nacional. O país foi dividido em norte, centro e sul, com a inicial de
cada zona a anteceder os algarismos. As chapas eram pretas com caracteres
brancos, com sequências como \plate{N-000}, \plate{C-000} e \plate{S-000}, ou
\plate{N-00000}, \plate{C-00000} e \plate{S-00000}.

\platefig{fmt-1911}{Matrícula da série de 1911: \plate{S-26384}.}

\begin{itemize}
  \item Nos automóveis de aluguer era colocado um \plate{A} no final do número de
        série. Ex.: \plate{C-4569-A}.
  \item Nas matrículas provisórias acrescentavam-se as letras \plate{WW}.
        Ex.: \plate{S-1703-WW}.
  \item Os automóveis do Exército usavam a sequência \plate{MG-000}, que se
        manteve em uso até à década de 1950.
\end{itemize}

\section{1918 | Zonas de registo dos Açores e da Madeira}

Foram atribuídas aos Açores e à Madeira, respetivamente, as séries \plate{A-000}
e \plate{M-000}. Ao contrário do que aconteceu no continente, este sistema
manteve-se em vigor nas ilhas até 1962.

\section{1937 | Série XX-00-00}
\subsection*{2.ª padronização do sistema de matrículas português}

Em algumas zonas o número de série já chegava aos cinco algarismos, o que
dificultava a rápida identificação da viatura. Foi nesta altura que começaram as
combinações que conhecemos, com duas letras e dois pares de algarismos.

A 1 de janeiro de 1937 entrou em vigor o 2.º sistema nacional de matriculação de
automóveis. Manteve-se a divisão por zonas, mas agora por grupos de letras:
\begin{itemize}
  \item de \plate{AA} a \plate{LZ} para a Zona de Lisboa;
  \item de \plate{MA} a \plate{TZ} para a Zona do Porto;
  \item de \plate{UA} a \plate{ZZ} para a Zona de Coimbra (antiga zona Centro).
\end{itemize}

As chapas continuaram a ser de fundo preto com caracteres a branco, com exceção
dos veículos diplomáticos ou em importação temporária, e eram de plástico, com
os caracteres em relevo.

\platefig{1937-pg0870}{Matrícula da série de 1937: \plate{PG-08-70}, da Zona do Porto.}

Não foi autorizado o uso das letras \plate{Q} e \plate{J}, por se confundirem
respetivamente com \plate{O} e com \plate{I}, e por razões de “decência” foi
interdito o uso dos grupos de letras \plate{CU} e \plate{FD}.

Alguns grupos de letras foram sendo reservados a organizações e entidades:\footnote{Segundo a Wikipédia (\emph{Matrículas de automóveis em Portugal}); a edição de 2023 não os referia.}
\begin{itemize}
  \item \plate{MG}, \plate{ME} e \plate{MX} para o Exército, \plate{AP} para a
        Marinha, \plate{AM} para a Força Aérea e \plate{EP} para a Administração
        Civil do Estado;
  \item \plate{CD}, \plate{CC} e \plate{FM}, em caracteres vermelhos sobre fundo
        branco, para o Corpo Diplomático, o Corpo Consular e outros funcionários
        em missão diplomática (até 1985, ver mais abaixo);
  \item \plate{TA}, \plate{TB} e \plate{TC}, em caracteres brancos sobre fundo
        vermelho, para os veículos em importação temporária das zonas de Lisboa,
        Porto e Coimbra.
\end{itemize}
A Guarda Nacional Republicana e a antiga Guarda Fiscal tinham sequências
próprias, começadas respetivamente por \plate{GNR} e \plate{GF}.

\subsection{Conversão das matrículas de 1911}

Os automóveis registados pelo sistema de 1911 foram obrigados a mudar para o
novo, através da atribuição de séries específicas de acordo com a matrícula
anterior. Os quatro algarismos finais do número antigo passaram a formar os dois
pares de algarismos.

\begin{itemize}
  \item \textbf{Lisboa:} \plate{S-000} e \plate{S-0000} passaram para \plate{AA-00-00};
        \plate{S-10000} para \plate{AB-00-00}; \plate{S-20000} para \plate{AC-00-00};
        \plate{S-30000} para \plate{AD-00-00}.
        Ex.: \plate{S-12345} passou para \plate{AB-23-45}.
  \item \textbf{Porto:} \plate{N-000} e \plate{N-0000} para \plate{MM-00-00};
        \plate{N-10000} para \plate{MN-00-00}.
        Ex.: \plate{N-1234} passou para \plate{MM-12-34}.
  \item \textbf{Coimbra:} todos os veículos mudaram para a série \plate{UU-00-00}.
        Ex.: \plate{C-123} passou para \plate{UU-01-23}.
\end{itemize}

\section{1962 | Adoção do sistema de matriculação nas ilhas}

Nos Açores e na Madeira o novo sistema só foi adotado em 1962. Os Açores foram
divididos em três zonas de matriculação, correspondentes aos Distritos Autónomos
de Ponta Delgada, Angra do Heroísmo e Horta, cada um com serviços próprios de
matrícula, e a Madeira constituía apenas uma zona, correspondente ao Distrito
Autónomo do Funchal. As letras reservadas a estas zonas foram:
\begin{itemize}
  \item \plate{AN} para Angra do Heroísmo;
  \item \plate{AR} e \plate{AS} para Ponta Delgada;
  \item \plate{HO} para a Horta;
  \item \plate{MA} e \plate{MD} para o Funchal.
\end{itemize}

As letras \plate{TD}, \plate{TG}, \plate{TH} e \plate{TF}, em caracteres brancos
sobre fundo vermelho, foram reservadas para os veículos em importação temporária
de Ponta Delgada, Angra do Heroísmo, Horta e Funchal, respetivamente.

\section{Anos 70 | Novos serviços regionais}

Na década de 1970 foram sendo estabelecidos novos serviços regionais e as
respetivas zonas de matrícula, com grupos de letras reservados:
\begin{itemize}
  \item \plate{EM} e \plate{EV} para Évora;
  \item \plate{ZA} e \plate{ZB} para Braga;
  \item \plate{ZC} e \plate{ZD} para Vila Real;
  \item \plate{ZE} e \plate{ZF} para Aveiro;
  \item \plate{ZG} e \plate{ZH} para a Guarda;
  \item \plate{ZI} e \plate{ZL} para Santarém;
  \item \plate{ZM} e \plate{ZN} para Setúbal;
  \item \plate{ZO} e \plate{ZP} para Faro.
\end{itemize}

O registo de Évora era considerado central a nível regional, pelo que lhe foi
também atribuído um código para veículos em importação temporária: \plate{TE}.
Como o registo por zonas foi descontinuado pouco tempo depois, a maioria destes
grupos de letras não chegou a ser usada como identificador de zona.

Com a independência das colónias portuguesas, quem quisesse podia importar o seu
carro para Portugal. Para o efeito foram atribuídas novas matrículas com as
iniciais \plate{VZ}, \plate{ZO}, \plate{ZE} e \plate{UZ}.\footnote{A atribuição de \plate{ZE} e \plate{UZ} a estas importações está por confirmar.}
\plate{VZ} era referido com humor como “Veio Zangado”.

\section{Anos 80 | Extensão do sistema de 1937}

No princípio da década de 1980, com o esgotamento das séries reservadas a algumas
zonas, acabou-se com a separação por zonas e os veículos passaram a ser
matriculados sequencialmente a nível nacional.

Mesmo assim, a série \plate{AA-00-00} aproximava-se rapidamente do fim. Para o
adiar, foi autorizada a utilização das “famosas” combinações \plate{CU} e
\plate{FD} e das letras \plate{Q} e \plate{J}.

Em 1985 as matrículas diplomáticas passaram a ter uma sequência diferente da
normal, nos formatos \plate{000-CD000}, \plate{000-CC000} e \plate{000-FM000}, em
que o primeiro grupo de algarismos identifica o país do utilizador do veículo e
o grupo de letras o seu estatuto diplomático (ver
secção~\ref{sec:diplomaticos}).

\platefig{diplomatic-fm}{Matrícula diplomática \plate{024-FM256}.}

\section{1992 | Série 00-00-XX}
\subsection*{Nova numeração e eurobanda}

Em 1992 as chapas mudaram de aparência e de numeração. Passaram a ser usadas
placas brancas refletoras com letras pretas e uma faixa azul do lado esquerdo
(eurobanda) com a identificação da União Europeia e a inicial do país, neste caso
\plate{P}. A fonte das letras é a DIN 1451, criada pelo Instituto Alemão de
Normalização em 1931. As novas chapas eram metálicas e Portugal foi o segundo
país a adotar a eurobanda.

A numeração passou a ter as letras no fim: \plate{00-00-AA}. As duas mudanças
não foram simultâneas: a nova chapa aplica-se às matrículas a partir de
1 de janeiro de 1992, mas a numeração \plate{00-00-AA} só começou a 1 de março
(Decreto-Lei n.º~54/2005, arts.~3.º e~5.º). Nos dois primeiros meses de 1992
foram, portanto, atribuídas matrículas \plate{AA-00-00} já em chapa branca.

Na nova sequência nenhum grupo de letras foi reservado a entidades especiais. A
única exceção foram as importações de 1997 (secção~\ref{sec:importados}).

\platefig{1992-7293fq}{Matrícula da série de 1992, sem barra amarela: \plate{72-93-FQ}.}

\section{1998 | Ano e mês do primeiro registo}

Mantendo a numeração de 1992, a aparência mudou: do lado direito passou a haver
uma faixa amarela com o ano e o mês do primeiro registo do veículo.

\platefig{1998-4572xq}{\plate{45-72-XQ}, com a barra amarela de junho de 2004.}

\section{2005 | Série 00-XX-00}

Esgotadas as combinações da numeração anterior, as matrículas passaram em 2005 a
ter o formato \plate{00-XX-00}.

\platefig{2005-99sa99}{Matrícula da série de 2005: \plate{99-SA-99}.}

\section{2020 | Série XX 00 XX}

Em 2020 mudaram tanto a numeração como a aparência das matrículas, para o
sistema usado atualmente. O Decreto-Lei n.º 2/2020, de 14 de janeiro, definiu o
novo número de matrícula como “dois grupos de duas letras e um grupo central de
dois algarismos”, um formato que o Decreto-Lei n.º~54/2005 já previa para quando
se esgotasse a série \plate{00-AA-00}. Segundo o IMT, aplica-se só a veículos
novos; o dono de um veículo de uma série anterior pode pedir uma chapa sem o mês
e o ano. A primeira matrícula, \plate{AA 01 AA}, foi emitida a 3 de março
de 2020 e atribuída a um veículo elétrico; a combinação \plate{AA 00 AA} ficou
reservada pelo IMT “para memória futura”.

\platefig{2020-aa16aa}{Matrícula da série de 2020: \plate{AA 16 AA}.}

Os traços foram substituídos por espaços, porque não haveria espaço suficiente
para números de matrícula com várias letras largas. A barra amarela com o ano e
o mês da primeira matrícula foi retirada para uniformizar as chapas com as dos
restantes países da União Europeia: Portugal era o único dos 27 Estados-Membros
com esta informação, que se podia confundir com a validade da matrícula que
alguns países têm nas suas chapas.

\subsection{A regra das vogais}

Com quatro letras, as combinações possíveis são variadas e algumas podem ser
menos apropriadas. Por isso, nas matrículas deste formato atribuídas pelo
sistema informático do IMT, não são usadas vogais na segunda nem na sexta
posição. Nas palavras do IMT:

\begin{quote}\small
“não são utilizadas vogais na segunda e sexta posição desses caracteres nas
matrículas do novo formato […] Tal não se aplica, quando qualquer dos grupos de
duas letras incluir duas vogais iguais.”
\end{quote}

Assim, \plate{AB 31 CE} não pode existir (vogal na 6.ª posição), mas
\plate{AA 12 BE} pode, porque o grupo \plate{AA} levanta a regra. O IMT refere
ainda que são evitadas “combinações de letras que possam formar determinadas
palavras”, sem publicar a lista.

\subsection{As letras K, W e Y}

Até 2020 as letras \plate{K}, \plate{W} e \plate{Y} não eram usadas nas
matrículas normais (à exceção do \plate{K} nas importações de 1997 e 1998). Com o
Acordo Ortográfico, o IMT anunciou que passariam a fazer parte da nova série;
não encontrámos confirmação de que já tenham sido emitidas.
